\section{Descrizione del problema}

Il dominio scelto per il nostro progetto è quello dell'\textbf{Arte e Cultura}. Il problema che intendiamo affrontare riguarda il progressivo allontanamento dei cittadini — in particolare dei giovani — dalla fruizione attiva dei luoghi culturali di prossimità, con un'attenzione specifica ai piccoli borghi e ai territori meno centrali.

I dati ISTAT 2023 confermano la gravità del fenomeno: solo una minoranza della popolazione italiana ha visitato musei, mostre o siti archeologici nell'ultimo anno. I giovani, pur essendo i fruitori più attivi, mostrano una partecipazione occasionale e non strutturata, mentre le fasce più anziane — portatrici di memoria storica — restano ai margini per barriere fisiche e digitali. A ciò si aggiunge un forte divario territoriale tra Centro-Nord e Mezzogiorno, e una partecipazione drasticamente inferiore nei piccoli comuni rispetto alle aree metropolitane, nonostante i borghi custodiscano patrimoni di altissimo valore identitario.

Il problema non è la mancanza di cultura, ma la sua \textbf{scarsa accessibilità, visibilità e capacità di coinvolgimento}. Le piattaforme esistenti (Italia.it, Google Arts \& Culture, Smartify) offrono esperienze prevalentemente passive, non orientate alla valorizzazione territoriale né alla partecipazione attiva dell'utente.

La nostra proposta consiste nell'evolvere \textbf{Culturia} — piattaforma partecipativa per la valorizzazione del patrimonio culturale — integrandovi funzionalità di \textbf{Realtà Aumentata (AR)} e \textbf{Intelligenza Artificiale (IA)}. Il sistema si articola su due direttrici tecnologiche complementari:

\begin{itemize}
    \item \textbf{Realtà Aumentata:} funge da strumento primario per l'\textbf{inclusività} e la \textbf{guida assistita} on-site. L'utente, inquadrando l'ambiente con il proprio smartphone, gode di un livello informativo sovrapposto alla realtà fisica che abbatte le barriere cognitive e linguistiche. Ricostruzioni visive, indicazioni spaziali sovraimpresse e approfondimenti multimodali rendono le visite autonome accessibili e altamente coinvolgenti anche nei luoghi sprovvisti di guide fisiche.

    \item \textbf{Intelligenza Artificiale:} viene integrata ad un duplice livello per operare in sinergia con l'AR. Per i visitatori, agisce come \textbf{motore di supporto alle guide in Realtà Aumentata}, adattando dinamicamente le informazioni mostrate e facendo da "cicerone" reattivo. Ancora più importante, l'IA assiste direttamente gli \textbf{enti territoriali e i gestori dei siti culturali} nella realizzazione dei tour stessi, abbassando la barriera tecnica necessaria per creare i contenuti AR tramite l'ausilio di interfacce generative.
\end{itemize}

Le due tecnologie operano in sinergia: l'IA adatta dinamicamente i contenuti AR al profilo dell'utente e può riconoscere automaticamente un sito culturale da una fotografia, arricchendolo con informazioni contestuali senza alcun input manuale. Queste funzionalità si integrano nell'ecosistema Culturia già progettato: il sistema di \textbf{gamification} (punti e badge per esperienze completate), il \textbf{diario delle visite} (ricordi multimediali arricchiti dai contenuti AR e dalle interazioni con l'IA) e la \textbf{validazione geolocalizzata} della presenza.

Il valore aggiunto di questa combinazione è duplice. Da un lato, rende fruibili contenuti altrimenti inaccessibili: un tempio in rovina può essere ``ricostruito'' virtualmente, un affresco deteriorato restituito nei colori originali. Dall'altro, supera la fruizione passiva offrendo un dialogo attivo e personalizzato con il patrimonio — un aspetto cruciale per i piccoli borghi privi di guide, segnaletica o personale, dove l'assistente IA diventa guida turistica e il tour AR supplisce all'assenza di infrastrutture.
